% !PRINTED 148
\weberpage{148}
%<<CONT
\[
 \nu_1,\nu_2,\nu_3\ldots\nu_n
\]
\gesperrt{eine abnehmende oder wenigstens niemals wachsende Zahlenreihe bilden}, oder dass die Differenzen
\[
 \nu_1-\nu_2,\nu_2-\nu_3,\ldots\nu_{n-1}-\nu_n
\]
alle positiv oder wenigstens nicht negativ sind.

Wir schliessen zweitens, dass in keinem Gliede der Function $S$ ein höherer Exponent als $\nu_1$ vorkommen kann; denn sonst würde auch ein Glied vorkommen, in dem $\alpha_1$ diesen höheren Exponenten hätte, und dies wäre gegen die Voraussetzung von höherer Ordnung als $A$. Es sind also die Glieder, die in einer symmetrischen Function überhaupt vorkommen können, von den Coëfficienten abgesehen, durch das höchste Glied vollkommen bestimmt, und die Anzahl der möglichen Glieder ist bei gegebenem höchsten Glied nur eine endliche.

Wir können die \gesperrt{Ordnung} der symmetrischen Function durch die Exponentenreihe ihres höchsten Gliedes
\[
 (\nu_1,\nu_2\ldots\nu_n)
\]
bezeichnen.

Von dem Fall abgesehen, in dem alle Exponenten Null sind, also die Function eine Constante ist, ist die möglichst niedrige Ordnung
\[
 (1,0,0\ldots0),
\]
der die einzige symmetrische Function
\[
 M(\alpha_1+\alpha_2+\cdots+\alpha_n)=-Ma_1
\]
entspricht.\begingroup\renewcommand{\thefootnote}{E2}\footnote[1]{[Quellenhinweis:] In diesen Beispielen bleibt ein beliebiger konstanter Summand unausgesprochen; ein solcher Summand ändert die angegebene Ordnung nicht. Die Formeln werden wie gedruckt beibehalten.}\endgroup{} Die nächst höhere Ordnung ist
\[
 (1,1,0\ldots0),
\]
der die symmetrische Function
\[
 Ma_1+M'a_2
\]
entspricht, und so erhalten wir in den $n$ ersten Ordnungen die Grundfunctionen selbst und ihre linearen Verbindungen, und keine anderen.

Die nächstfolgende Ordnung ist charakterisirt durch
\[
 (2,0,0\ldots0)
\]
und enthält die linearen Verbindungen der Grundfunctionen mit der Summe der Quadrate.

In allen diesen Fällen ist die Darstellbarkeit der symmetrischen Functionen bereits bewiesen, und wir werden also jetzt den allgemeinen Beweis dadurch führen, dass wir die Darstellung
%CONT>>
